\chapter{전시용 메뉴 카드: \texorpdfstring{\code{tools/menu\_cards}}{tools/menu\_cards}}\label{ch:32-menu-cards}
\begingroup\setlength{\emergencystretch}{3em}% long identifiers: allow looser lines in this chapter only

전시의 실습은 카드 한 장에서 시작한다. 관람객은 \term{메뉴 카드}{menu card}를 고르고, 카드에 인쇄된 대로 기계를
배선하고, 돌리고, 정지를 읽고, 카드 뒷면을 뒤집어 답을 맞춰 본다. 기어식, 릴레이식, FPGA의 세 기계에 \emph{같은
카드}를 걸어 같은 정지를 보여 주는 것이 전시의 핵심이다(\code{hardware/common/design-brief.md}). 이 장은 그 카드를
만드는 프로그램 \code{tools/menu\_cards}와, 그것이 만든 기록 \code{hardware/common/menu-cards/*.json}을 읽는다.

카드에는 암호문, 크립, 메뉴 그림, 기계 설정이 있다. 그 뒤에는 ``이 카드로 돌리면 정지가 어디서 몇 번 나오고, 체킹하면
어떤 평문이 나온다''는 약속이 있다. 프로그램이 하는 일은 그 약속이 참인 카드를 \emph{찾는} 것, 그리고 그것이 참임을
\emph{다시 확인하는} 것이다.

\section{카드의 모양: 두 가지 시나리오}\label{sec:mc-shapes}

\begin{lstlisting}[language=Python,caption={\code{scenarios.py}: 기본 카드 여섯 장의 요구 조건.},label={lst:mc-scenarios}]
DEFAULT_SET: list[tuple[Spec, int]] = [
    (Spec("slice-6-a", "slice", 6, 4, max_span=12, false_min=0, false_max=0), 1),
    (Spec("slice-6-b", "slice", 6, 4, max_span=12, false_min=0, false_max=0), 2),
    (Spec("slice-6-c", "slice", 6, 4, max_span=12, false_min=1, false_max=1), 3),
    (Spec("slice-6-d", "slice", 6, 4, max_span=12, false_min=1, false_max=2), 4),
    (Spec("full-12-a", "full", 12, 10, max_span=16, false_min=0, false_max=0), 11),
    (Spec("full-12-b", "full", 12, 10, max_span=16, false_min=1, false_max=1), 12),
]
\end{lstlisting}

\code{Spec}은 카드가 갖춰야 할 조건이다. 구동 방식, 선의 수(곧 스크램블러 유닛의 수), 글자의 수, 선이 들어갈 크립
구간의 최대 길이, 그 구동에서 나올 거짓 정지 수의 범위. 마지막 수는 탐색의 씨앗이다(\ref{sec:mc-seed}절).
루프의 수는 따로 적지 않는다. 한 덩어리로 이어진 메뉴에서 루프는 (선 $-$ 글자 $+ 1$)이므로(\ref{ch:rust-menu}장),
두 시나리오 모두 루프가 셋이다.

\begin{itemize}
\item \textbf{조각 카드}(\code{slice}, 유닛 6개). 글자 4개, 선 6개, 루프 3개. 느린 드럼을 정답 위치의 글자에 세워 두고
나머지 두 드럼만 676 위치를 돈다(\ref{sec:s-slice}절). 관람객의 자유 실습용이다. 릴레이식으로 43초, 실습 전체로
7--8분이다.
\item \textbf{전체 카드}(\code{full}, 유닛 12개). 글자 10개, 선 12개, 루프 3개. 카드에 적힌 로터 순서 하나의 17,576
위치를 모두 돈다. 해설자의 시연용이다.
\end{itemize}

\Why 글자 4개에 선 6개는 특별한 모양이다. 글자 네 개 사이에 서로 다른 짝은 $\binom{4}{2} = 6$개뿐이므로, 같은 짝을
두 번 쓰지 않는다면(생성기는 그렇게 요구한다) 모든 글자 쌍이 선 하나로 이어진 \term{완전 그래프}{complete graph}
$K_4$가 될 수밖에 없다(그림~\ref{fig:g-mc-k4}). 스크램블러 여섯 개로 얻을 수 있는 가장 촘촘한 메뉴다. 루프가 많을수록
거짓 정지가 적고(\ref{ch:rust-menu}장), 전시의 짧은 실습에서 관람객이 체킹할 수 있는 정지는 한두 개이므로 이 모양을
택했다.

\begin{figure}[ht]
\centering
\begin{tikzpicture}[x=1cm,y=1cm,el/.style={circle,draw=corange,fill=white,text=corange,inner sep=0.5pt,minimum size=4.6mm,font=\scriptsize\bfseries}]
  \node[ltr,minimum size=8mm,font=\ttfamily\bfseries,draw=corange,line width=1.3pt] (D) at (0,1.6) {D};
  \node[ltr,minimum size=8mm,font=\ttfamily\bfseries] (N) at (1.6,0) {N};
  \node[ltr,minimum size=8mm,font=\ttfamily\bfseries] (U) at (0,-1.6) {U};
  \node[ltr,minimum size=8mm,font=\ttfamily\bfseries] (Z) at (-1.6,0) {Z};
  \draw (N) -- node[el] {3} (Z);
  \draw (Z) -- node[el] {4} (U);
  \draw (U) -- node[el] {7} (N);
  \draw (D) -- node[el] {9} (Z);
  \draw (U) -- node[el,pos=0.3] {11} (D);
  \draw (N) -- node[el] {12} (D);
  \node[note,anchor=west,text=corange] at (0.55,1.95) {테스트 글자};
  \begin{scope}[xshift=4.0cm,yshift=0.3cm]
    \matrix[matrix of nodes,nodes={font=\ttfamily\footnotesize,inner sep=1.2pt,minimum width=4.4mm,anchor=center},
            column sep=0pt,row sep=1pt] (m) {
      |[font=\scriptsize]| 링크 & 1 & 2 & |[text=corange,font=\bfseries\ttfamily\footnotesize]| 3 & |[text=corange,font=\bfseries\ttfamily\footnotesize]| 4 & 5 & 6 & |[text=corange,font=\bfseries\ttfamily\footnotesize]| 7 & 8 & |[text=corange,font=\bfseries\ttfamily\footnotesize]| 9 & 10 & |[text=corange,font=\bfseries\ttfamily\footnotesize]| 11 & |[text=corange,font=\bfseries\ttfamily\footnotesize]| 12 \\
      |[font=\scriptsize]| 크립 & B & E & N & Z & I & N & U & N & D & M & U & N \\
      |[font=\scriptsize]| 암호 & W & N & Z & U & J & S & N & H & Z & N & D & D \\
    };
    \foreach \c in {4,5,8,10,12,13} { \draw[corange,rounded corners=1pt] (m-1-\c.north west) rectangle (m-3-\c.south east); }
    \node[note,anchor=north west,align=left] at (m-3-1.south west) {크립 \texttt{BENZINUNDMUN}, 전문의 90--101번째 글자.\\
      색 칸이 카드의 스크램블러 6개다.};
  \end{scope}
\end{tikzpicture}
\caption{카드 \code{slice-6-a}의 메뉴. 글자 4개 사이의 짝 6개가 모두 선이다($K_4$). 선의 번호는 크립 위치
(링크 번호)이고, 오른쪽은 카드 앞면의 정렬 띠를 옮긴 것이다. 값은 \code{hardware/common/menu-cards/slice-6-a.json}에서
가져왔다.}
\label{fig:g-mc-k4}
\end{figure}

\section{자연스러운 크립으로는 안 되는 이유}\label{sec:mc-why-engineer}

보통의 크립에서 메뉴를 만들면 글자가 많고 루프가 적다. 노트북의 예에서 16글자 크립은 글자 18개, 루프 1개의 메뉴를
주었다(\ref{ch:15-python-api}장). 암호문의 글자는 사실상 무작위이므로, 선 6개의 암호문 글자가 모두 정해진 네 글자
안에 떨어지는 일은 거의 없다. 크립을 바꿔 가며 기다려서는 $K_4$를 얻을 수 없다.

그래서 생성기는 반대로 간다. 메뉴를 먼저 정하고 그 메뉴를 주는 \emph{키}를 만든다. 특히 플러그보드를 고른다.

\Why 에니그마로 크립 위치 $k$의 평문 글자 $p_k$를 치면 전류는 플러그보드 $S$, 스크램블러 $\sigma_k$, 다시 $S$를
지난다(\ref{sec:bombe-one-page}절).
\[
  c_k = S\bigl(\sigma_k(S(p_k))\bigr)
\]
로터 순서, 링, 시작 위치를 먼저 정하면 $\sigma_k$는 모두 계산할 수 있다(플러그보드 없는 에니그마). 남은 자유는
$S$뿐이다. $S$를 잘 고르면 $c_k$가 원하는 글자 안에 떨어지게 할 수 있다. 플러그보드는 1940년대 키의 가장 큰
부분이었으므로, 그것을 마음대로 고를 수 있다면 메뉴의 모양을 상당히 고를 수 있다.

\General 시험 문제를 만드는 선생님은 답이 깔끔한 수가 되도록 문제의 숫자를 거꾸로 정한다. ``답이 $x=3$이 되게 하려면
계수를 무엇으로 할까''. 여기서도 원하는 결과(메뉴의 모양)에서 출발해 입력(플러그보드)을 정한다.

\section{플러그보드를 찾는 깊이 우선 탐색}\label{sec:mc-dfs}

$S$는 아무 표나 될 수 없다. 플러그보드는 글자 쌍을 맞바꾸므로 $S(x)=y$이면 $S(y)=x$여야 한다(자기 자신과 짝인 글자는
플러그가 없는 것이다). 이 제약을 지키면서 선 하나씩 $S$의 값을 정해 나간다.

\begin{lstlisting}[language=Python,caption={\code{engine.py}: \code{solve\_plugboard}의 핵심(일부).},label={lst:mc-dfs}]
def rec(i: int, mask: int) -> bool:
    ...
    if i == E:
        if popcount(mask) == n_letters and connected():
            found.append(list(S))
            return True
        return False
    p = plains[i]
    new_p = S[p] < 0
    opts_a = [p] + [x for x in order if S[x] < 0 and x != p] if new_p else [S[p]]
    for a in opts_a:
        if new_p:
            S[p] = a
            S[a] = p
        y = sigmas[i][a]
        if S[y] >= 0:
            opts_c = [S[y]]
        else:
            inm = [x for x in order if mask >> x & 1 and S[x] < 0 and x != y]
            oth = [] if popcount(mask | 1 << p) >= n_letters else [...]
            opts_c = [y] + inm + oth
        for c in opts_c:
            ...                      # skip c == p, repeated pairs, too many letters
            S[y] = c; S[c] = y       # (if S[y] was still free)
            if rec(i + 1, mask | 1 << p | 1 << c):
                return True
            ...                      # undo
        ...                          # undo S[p]
    return False
\end{lstlisting}

$i$번째 선에서 정할 것은 둘이다.

\begin{enumerate}
\item $a = S(p_i)$. 평문 글자 $p_i$의 짝이 아직 없으면 자기 자신(플러그 없음)부터 비어 있는 글자를 차례로 시험한다.
이미 정해져 있으면 그 값 하나뿐이다.
\item $c = S(y)$, 여기서 $y = \sigma_i(a)$. 이것이 그 선의 암호문 글자다. $y$의 짝이 이미 정해져 있으면 선택의 여지가
없다. 비어 있으면 $y$ 자신, 이미 메뉴에 들어온 글자, 그리고 글자 수에 여유가 있을 때만 새 글자를 시험한다.
\end{enumerate}

조건을 어기는 가지는 그 자리에서 버린다. 암호문 글자가 평문 글자와 같으면 안 되고, 같은 글자 쌍이 두 번 나오면 안 되며
(단순 그래프), 글자가 \code{n\_letters}개를 넘으면 안 된다. 끝까지 내려가면 글자가 정확히 그 수이고 메뉴가 한 덩어리인지
본다. 실패하면 마지막에 정한 값을 지우고(``undo'') 다음 후보로 간다(그림~\ref{fig:g-mc-tree}).

\begin{figure}[ht]
\centering
\begin{tikzpicture}[x=1cm,y=-1cm,
    st/.style={draw,rounded corners=1.5pt,fill=white,font=\scriptsize,inner sep=2pt,align=center},
    bad/.style={st,fill=credbg,draw=cred},
    good/.style={st,fill=cgreenbg,draw=cgreen}]
  \node[st] (r) at (0,0) {선 1: $S(p_1)$?};
  \node[st] (a1) at (-3.5,1.2) {$S(p_1)=p_1$};
  \node[st] (a2) at (0,1.2) {$S(p_1)=Q$};
  \node[note] at (3.0,1.2) {\ldots};
  \node[st] (c1) at (-5.2,2.4) {$S(y_1)=y_1$};
  \node[bad] (c2) at (-2.4,2.4) {$c_1 = p_1$\\버림};
  \node[st] (d1) at (-5.2,3.6) {선 2: \ldots};
  \node[bad] (d2) at (-6.8,4.8) {글자 5개째\\버림};
  \node[bad] (d3) at (-3.9,4.8) {같은 쌍\\버림};
  \node[st] (c3) at (0,2.4) {$S(y_1)=N$};
  \node[note] at (0,3.6) {\ldots};
  \node[good] (ok) at (0,4.8) {선 6까지: 글자 4개,\\한 덩어리 $\Rightarrow$ 찾음};
  \draw[->] (r) -- (a1); \draw[->] (r) -- (a2); \draw[->] (a1) -- (c1); \draw[->] (a1) -- (c2);
  \draw[->] (c1) -- (d1); \draw[->] (d1) -- (d2); \draw[->] (d1) -- (d3);
  \draw[->] (a2) -- (c3); \draw[->] (c3) -- (0,3.4); \draw[->] (0,3.8) -- (ok);
  \draw[->,dashed,cgray] (d3.east) to[bend right=20] node[note,above,pos=0.6] {되돌아감} (c3.west);
\end{tikzpicture}
\caption{플러그보드의 깊이 우선 탐색(모양만 그린 것). 선마다 $S(p_i)$와 $S(y_i)$를 정하고, 조건을 어기는 가지는
그 자리에서 버리고 되돌아간다. 마디를 4,000개 넘게 보면 이 키를 포기한다.}
\label{fig:g-mc-tree}
\end{figure}

\Why 이 방법은 \term{되추적}{backtracking}이다. 모든 플러그보드를 늘어놓고 하나씩 검사하면 $10^{14}$ 단위의 수를 봐야
하지만, 선 하나를 정할 때마다 조건을 검사해 틀린 가지를 일찍 자르면 볼 것이 크게 준다. 탐색은 마디 4,000개를 넘으면
멈추고(\code{budget}) 실패로 돌아간다. 이 키로는 잘 안 되는 것이니 다른 키를 뽑는 편이 빠르다.

\General 스도쿠를 손으로 풀 때 빈칸에 숫자를 하나 써 보고, 같은 줄에 겹치는 숫자가 생기면 지우고 다음 숫자를 써 보는
것이 되추적이다. 여덟 개의 퀸을 서로 공격하지 않게 놓는 문제도 같은 방법으로 푼다.

\Ours 후보를 시험하는 순서(\code{order})는 무작위로 섞는다. 같은 크립에서도 매번 다른 플러그보드를 찾게 하기 위해서다.
찾은 $S$가 정한 짝은 대개 10쌍보다 적으므로, \code{plugboard\_string}이 빈 글자로 짝을 더 채워 실제 운용처럼 10쌍을
만든다. 메뉴에 쓰이지 않는 글자의 짝이므로 메뉴의 모양에는 영향이 없다.

\section{후보 하나 만들기}\label{sec:mc-candidate}

플러그보드 탐색은 \code{build\_candidate}의 한 단계다. 후보 하나는 이렇게 만든다.

\begin{enumerate}
\item \textbf{전문.} \code{texts.py}의 독일군 전문 여덟 개 가운데 하나를 고른다. 대문자만 쓰고, 움라우트는
\texttt{AE}, \texttt{OE}, \texttt{UE}로 풀어 쓰고, 낱말 사이에 \texttt{X}를 넣고, 사람 이름은 넣지 않는다는 실제 통신의
관례를 따른다. 전문마다 크립 후보가 될 문구가 셋씩 적혀 있다(\texttt{WETTERVORHERSAGE}, \texttt{BENZINUNDMUNITION} 등).
\item \textbf{키.} 로터 I--V 가운데 셋, 반사판 B, 링과 시작 위치는 무작위로 뽑는다(Enigma~I).
\item \textbf{크립 구간.} 문구 안에서 길이 \code{max\_span}(조각 12, 전체 16)의 구간을 고른다. 그 구간에서 왼쪽과
가운데 로터의 창 글자가 변하지 않아야 한다. 곧 크립 안에서 가운데 로터가 넘어가지 않아야 한다.
\item \textbf{선의 자리.} 선으로 쓸 크립 위치를 고른다(\code{pick\_positions}). 선의 평문 글자는 많아야 메뉴의 글자
수만큼이어야 하므로, 크립에 자주 나오는 글자가 있는 자리를 고른다. \texttt{BENZINUNDMUN}에는 \texttt{N}이 넷,
\texttt{U}가 둘 있다.
\item \textbf{플러그보드.} 앞 절의 탐색으로 $S$를 찾는다.
\item \textbf{실제 소프트웨어로 다시 만들기.} 찾은 키로 \emph{실제} \code{Enigma}가 전문 전체를 암호화하고, \emph{실제}
\code{Menu.from\_crib}이 메뉴를 만들고, \code{subset}이 그 자리들의 부분 메뉴를 만든다(\ref{sec:s-subset}절). 부분 메뉴의
선, 글자, 루프 수가 요구와 정확히 같을 때만 후보로 남긴다.
\end{enumerate}

\Why 3번의 ``넘어감 없음''은 실물 기계를 위한 조건이다. 봄베는 크립 안에서 가운데 로터가 움직이지 않는다고 가정하고,
그렇지 않으면 맞는 위치에서도 정지하지 않는다(\ref{ch:rust-bombe}장). 소프트웨어는 넘어감 가설로 이것을 다룰 수 있지만
실물 두 대에는 그런 장치(실제 봄베의 ``hoppity'')가 없다. 카드가 그 문제를 아예 피하게 만든다.

6번이 이 생성기의 원칙이다. 탐색은 자기 나름의 계산($\sigma_k$의 표)으로 플러그보드를 찾지만, 그 계산을 믿고 끝내지
않는다. 키가 정해지면 그다음은 모두 B부에서 읽은 기준 코드가 다시 한다. 생성기에 계산 실수가 있으면 6번에서 모양이
어긋나 후보가 버려질 뿐, 틀린 카드가 나오지는 않는다.

\section{후보를 평가하기}\label{sec:mc-evaluate}

모양이 맞는 후보를 봄베와 체킹에 걸어 본다(\code{evaluate}).

\begin{lstlisting}[language=Python,caption={\code{engine.py}: 조각 구동과 정답 위치.},label={lst:mc-stops}]
def stops_for(menu: Menu, order: WheelOrder, mode: str, true_right: str | None):
    """Run the software bombe over one wheel order; in slice mode only the 676 positions with the
    right (slow) drum fixed at ``true_right``."""
    bombe = Bombe(MODEL).plug_up(menu)
    stops = bombe.run(order, fixed_right=true_right if mode == "slice" else None)
    return bombe, stops

def true_stop_letters(c: Candidate) -> str:
    """Drum letters of the true stop = core positions of the rotors at the first crib letter."""
    rings = c.key["rings"]
    w = c.states[0]
    return "".join(L((N(w[i]) - N(rings[i])) % 26) for i in range(3))
\end{lstlisting}

정답 위치는 키에서 바로 계산된다. 크립의 첫 글자를 칠 때의 창 글자에서 링 설정을 뺀 코어 위치다
(\ref{ch:rust-check}장의 ``코어 위치''). 조각 카드는 그 셋째 글자, 곧 오른쪽(느린) 드럼의 글자에 드럼을 세우고 돈다.
그런 다음 차례로 확인한다.

\begin{itemize}
\item 정답 위치의 정지가 정확히 하나 있다.
\item 거짓 정지의 수가 \code{false\_min}--\code{false\_max} 안에 있다.
\item 모든 정지를 체킹한다. 정답 위치에서는 전문 전체가 정확히 복원되어야 하고, 거짓 정지는 하나도 체킹을 통과하면 안
된다.
\end{itemize}

마지막 조건은 관람객을 위한 것이다. 거짓 정지가 체킹까지 통과하면 관람객은 틀린 답을 정답으로 믿게 된다. 그런 카드는
쓰지 않는다. 기본 카드에는 거짓 정지가 0개, 1개, 2개인 것이 섞여 있다(코드~\ref{lst:mc-scenarios}). 어떤 카드에서는
관람객이 정답만 만나고, 어떤 카드에서는 거짓 정지를 만나 체킹으로 그것을 버리게 된다.

\paragraph{참고로 세는 것.} 평가는 카드의 구동 말고도 두 수를 더 센다. 같은 메뉴로 정답 로터 순서의 전체
17,576 위치를 돌았을 때의 거짓 정지와, \code{run("all")}로 모든 로터 순서를 돌았을 때의 거짓 정지다. Enigma~I에서
\code{"all"}은 반사판 B와 C를 합쳐 120가지다(\ref{sec:s-count}절). 기록의 필드
\code{false\_stops\_all\_120\_orders\_full\_run}이 이 120가지를 모두 돈 결과다. 관람객이 하지 않는 구동이지만, 조각과 카드가 얼마나 일을 줄여 주는지를 보여 주는 수다.

\begin{table}[ht]
\centering
\caption{기본 카드 여섯 장(\code{hardware/common/menu-cards/*.json}). 거짓 정지는 카드의 구동 / 정답 로터 순서의 전체 구동 /
\code{run("all")}의 순서다.}
\label{tab:g-mc-cards}
\small\setlength{\tabcolsep}{4pt}
\begin{tabular}{lllllrr}
\toprule
카드 & 로터 순서 & 크립 & 메뉴 글자 & 정답 위치 & 거짓 정지 & 시도 번호 \\
\midrule
\code{slice-6-a} & B III-II-IV & \texttt{BENZINUNDMUN} & \texttt{DNUZ} & \texttt{KMX} & 0 / 12 / 1,881 & 2,232 \\
\code{slice-6-b} & B IV-III-I & \texttt{LAERUNGMELDE} & \texttt{DELN} & \texttt{NDZ} & 0 / 17 / 1,903 & 14,892 \\
\code{slice-6-c} & B I-II-IV & \texttt{NZINUNDMUNIT} & \texttt{IMNU} & \texttt{RZP} & 1 / 16 / 1,883 & 6,778 \\
\code{slice-6-d} & B IV-V-III & \texttt{RWUNDETEWERD} & \texttt{DERW} & \texttt{NWC} & 2 / 21 / 1,911 & 17,525 \\
\code{full-12-a} & B I-II-V & \texttt{NDERENEREIGNISSE} & \texttt{DEGINORSXZ} & \texttt{SWJ} & 0 / 0 / 46 & 133 \\
\code{full-12-b} & B II-I-V & \texttt{KEINEBESONDERE} & \texttt{BDEGIKNRSW} & \texttt{VFI} & 1 / 1 / 53 & 568 \\
\bottomrule
\end{tabular}
\end{table}

표~\ref{tab:g-mc-cards}에서 조각 카드의 숫자를 보면 조각 구동의 뜻이 보인다. 같은 $K_4$ 메뉴로 정답 로터 순서를
끝까지 돌면 거짓 정지가 12--21개 나온다. 설계 기준서가 이 모양의 메뉴에 대해 측정에서 잡은 평균 17.8개와 같은
규모다. 관람객이 7분 안에 체킹하기에는 너무 많다. 느린 드럼을 세우면 그 $1/26$만 남고, 생성기가 그중 0--2개인 카드를
고른다. 글자 10개짜리 전체 카드는 루프는 같아도 글자가 많아서 로터 순서 전체에서도 거짓 정지가 46, 53개뿐이다.

\section{같은 씨앗, 같은 카드}\label{sec:mc-seed}

\begin{lstlisting}[language=Python,caption={\code{engine.py}: 시도 하나는 씨앗과 번호로 정해진다.},label={lst:mc-attempt}]
def attempt(spec: Spec, seed: int, index: int, msg_ids=None):
    """Attempt number ``index`` of a deterministic search; returns (candidate, evaluation) or None."""
    rng = random.Random(seed * 1_000_003 + index)
    pool = [m for m in texts.MESSAGES if not msg_ids or m[0] in msg_ids]
    c = build_candidate(spec, rng, msg=rng.choice(pool))
    ...
\end{lstlisting}

시도 하나에 쓰는 난수는 모두 (씨앗, 시도 번호)에서 만든 난수 생성기 하나에서 나온다. 그러므로 시도 번호를 알면 그
시도를 언제든 똑같이 되풀이할 수 있다. 탐색은 시도 0, 1, 2, \ldots{}를 차례로 해서 처음 성공한 번호를 기록한다.
여러 프로세스로 나눠 돌릴 때도 마찬가지다. 각 프로세스가 연속된 번호 20개씩을 맡고, 한 차례에 성공한 번호들
가운데 가장 작은 것을 고른다. 그보다 작은 번호는 앞 차례들에서 모두 실패했으므로, 결과는 일꾼 수와 무관하게
``가장 작은 성공 번호''다.

\Why 카드는 전시장에 인쇄되어 오래 남는다. 누가 언제 다시 만들어도 같은 카드가 나와야 카드에 문제가 생겼을 때 원인을
따라갈 수 있다. 기록에 씨앗과 시도 번호를 적어 두므로, 시험 \code{test\_search\_is\_deterministic}이 카드마다 그
번호의 시도를 다시 해서 같은 암호문, 같은 정답 위치, 같은 거짓 정지 수가 나오는지 본다. 표~\ref{tab:g-mc-cards}의
시도 번호가 그것이다. 조각 카드는 수천에서 만 칠천 번째 시도에서야 조건을 다 갖춘 키가 나왔다.

\General 컴퓨터 게임에서 같은 \emph{시드}를 넣으면 같은 지도가 만들어지는 것과 같다. 지도 전체를 저장하는 대신 시드
하나만 알려 주면 다른 사람도 같은 세계를 얻는다.

\section{기록과 독립적인 재확인}\label{sec:mc-record}

찾은 카드는 JSON 기록 하나로 저장된다(\code{record.make\_record}). 암호문(다섯 글자씩 묶은 것 포함), 평문, 크립과 그
위치, 메뉴(유닛마다 평문·암호문 글자, 링크 번호, 드럼 오프셋), 로터 순서, 정답 위치(코어 글자, 에니그마 창 글자),
예상 정지 목록과 거짓 정지 수, 키, 체킹이 복원한 키가 들어 있다.

\paragraph{유닛의 드럼 오프셋.} 크립 위치 $k$의 유닛은 오른쪽 드럼이 $k$칸 앞서고 왼쪽, 가운데 드럼은 같다
(\ref{ch:rust-bombe}장). 기록은 이것을 \code{"drum\_offset\_letters": "AA"+}($k$번째 글자)로 적는다. \code{slice-6-a}의
링크 3 유닛은 $k=2$이므로 \texttt{AAC}다. 기록을 만들 때 \code{assert}로 한 가지를 더 확인한다. 정답 위치의 코어 글자에
링 설정을 더하면 에니그마의 창 글자가 되어야 한다. 이 프로젝트에서 위치를 나타내는 방법이 셋(코어, 창 글자, 드럼
오프셋)이므로 셋 사이의 환산이 맞는지 기록마다 본다.

\paragraph{다시 확인하기.} \code{verify\_record}는 기록의 \emph{필드만} 가지고 처음부터 다시 계산한다.

\begin{enumerate}
\item 키와 평문으로 암호화하면 기록의 암호문이 나온다.
\item 크립이 평문의 그 위치에 있고, 크립 글자와 그 아래 암호문 글자가 같은 자리가 없다.
\item 크립 구간에서 왼쪽, 가운데 로터가 움직이지 않는다.
\item 기록의 크립 위치들로 만든 부분 메뉴의 선, 글자, 루프 수, 테스트 글자, 입력 선이 기록과 같다.
\item 봄베(조각이면 기록의 느린 드럼 글자로)를 돌린 정지의 집합이 기록과 같고, 정답 위치가 그 안에 있고, 거짓 정지 수가
같다.
\item 창 글자 $-$ 링이 정답 위치다.
\item 체킹하면 정답 위치에서 평문이 복원되고, 거짓 정지는 통과하지 않는다.
\end{enumerate}

\Why 생성기가 스스로 확인한 것을 왜 또 확인하는가. \code{verify\_record}는 생성기의 중간 계산(플러그보드 탐색, 후보
객체)을 하나도 쓰지 않는다. 기록 파일과 Rust 코어만 있으면 된다. 그래서 생성기가 틀려도, 누가 기록 파일을 손으로
고쳐도, 나중에 코어가 바뀌어도 이 확인이 잡아낸다. 시험 \code{test\_tampered\_record\_is\_rejected}는 기록의
플러그보드를 지우거나 루프 수를 하나 늘리면 \code{verify\_record}가 문제를 보고하는지 본다. 명령
\code{python -m menu\_cards verify ../hardware/common/menu-cards/\{full,slice\}-*.json}으로(목록 파일 \code{index.json}은 빼고) 언제든 다시 확인할 수 있다.

\General 회계 장부를 쓴 사람과 감사하는 사람을 나누는 것과 같다. 감사인은 장부를 쓴 사람의 계산 과정을 믿지 않고,
영수증(기록)만 가지고 다시 더한다.

\Ours 기록에 적힌 정지 하나하나는 실물 기계의 시험에도 쓰인다. 설계 기준서는 개관 전에 기준 메뉴 카드를 걸어
알려진 정지 위치와 스테커 후보가 나오는지 확인하도록 적고 있다. 실제 봄베의 운용에서 ``check stop''으로 기계를
점검하던 것과 같은 생각이다. 카드는 관람객의 교재이자 기계의 시험지다.

\section{카드를 그리기}\label{sec:mc-card}

\code{card.py}는 기록에서 A5 크기 카드 한 장의 LaTeX 원고를 만들고 xelatex로 PDF를 만든다. 앞면과 뒷면이다
(그림~\ref{fig:g-mc-layout}).

\begin{figure}[ht]
\centering
\begin{tikzpicture}[x=1cm,y=1cm,
    sec/.style={draw=cgray!60,fill=white,minimum width=50mm,align=left,anchor=north west,font=\scriptsize,inner sep=2pt,text width=48mm}]
  \foreach \xo/\ttl in {0/앞면, 6.2/뒷면} {
    \draw[fill=corangebg!35,draw=cgray] (\xo,0) rectangle (\xo+5.6,7.6);
    \node[font=\small\bfseries,anchor=south] at (\xo+2.8,7.65) {\ttl};
  }
  \node[sec] at (0.3,7.3) {\textbf{메뉴 카드} \hfill \texttt{slice-6-a}\\스크램블러 6 $\cdot$ 글자 4 $\cdot$ 루프 3};
  \node[sec] at (0.3,6.45) {1. 전문: 다섯 글자 묶음};
  \node[sec] at (0.3,5.75) {2. 크립: 정렬 띠(링크, 크립, 암호)};
  \node[sec,minimum height=24mm] at (0.3,5.05) {3. 메뉴 그래프\\(테스트 글자 굵은 테두리)};
  \node[sec] at (0.3,2.4) {4. 기계 설정: 반사판, 로터,\\느린 드럼 \texttt{X}에 고정, 676 위치};
  \node[sec] at (6.5,7.3) {\textbf{배선 안내}};
  \node[sec,minimum height=12mm] at (6.5,6.7) {유닛 $\mid$ 링크 $\mid$ IN $\mid$ OUT $\mid$ 드럼 오프셋};
  \node[sec] at (6.5,5.3) {테스트 레지스터: 글자 \texttt{D}, 입력 선 \texttt{A}};
  \node[sec,minimum height=12mm] at (6.5,4.6) {정지 기록(빈 칸): 지시 드럼, 꺼진 램프, 체킹 결과};
  \node[sec,minimum height=17mm,rotate=180,anchor=south east] at (6.5,1.95) {정답: 정지, 키, 해독문};
  \node[note,anchor=west] at (11.9,1.1) {180°\\돌려 인쇄};
\end{tikzpicture}
\caption{카드의 배치(\code{card.py}의 \code{front}와 \code{back}을 줄여 그린 것). 정답은 뒷면 아래쪽에 거꾸로
인쇄된다.}
\label{fig:g-mc-layout}
\end{figure}

\paragraph{메뉴 그래프의 배치.} 글자를 원 위에 늘어놓고 선을 곧게 긋는다. 원 위의 순서에 따라 선이 서로 몇 번
엇갈리는지가 달라지므로, \code{circle\_order}가 엇갈림이 적은 순서를 찾는다. 무작위 순서에서 시작해 두 글자의 자리를
바꿔 보고 엇갈림이 줄면 받아들이는 일을 더 줄지 않을 때까지 되풀이하고, 이것을 시작점을 바꿔 60번 한다.
\ref{ch:rust-check}장의 언덕 오르기와 같은 방법이다. 엇갈림이 0이 되면 일찍 멈춘다. $K_4$는 어떻게 놓아도 엇갈림이
하나 남는다. 원 위의 네 점을 모두 이으면 두 대각선이 반드시 가운데에서 만나기 때문이다. 그 두 대각선의 번호가 가운데에서
겹치지 않도록, 번호를 놓을 자리를 선의 가운데부터 시작해 다른 번호와 일정 거리 이상 떨어진 곳을 찾아 옮긴다.
그림~\ref{fig:g-mc-k4}에서 두 대각선은 링크 3과 11이고, 11의 번호가 가운데를 비켜 있는 것이 그 결과를 흉내 낸
것이다.

\paragraph{정답을 거꾸로.} 뒷면 아래의 정답(정지 목록, 키, 해독문)은 \code{\textbackslash rotatebox\{180\}}으로
뒤집어 인쇄한다. 카드를 바로 들고 있으면 읽기 어렵고, 다 푼 뒤 카드를 돌리면 읽힌다. 퍼즐 잡지의 정답란과 같은
배려다.

\paragraph{PDF는 저장소 밖에.} 기록(JSON)은 저장소에 두지만, PDF와 LaTeX의 중간 파일은 사용자의 로컬 데이터 폴더
아래 \code{christopher/menu-cards}에 만든다(바꾸려면 \code{-{}-pdf-dir}). PDF는 기록과 \code{card.py}에서 언제든 다시
만들 수 있으므로 기준이 아니다. 기준은 기록이다.

\section{시험}\label{sec:mc-tests}

\code{tests/test\_menu\_cards.py}가 확인하는 것을 정리한다.

\begin{center}
\small
\begin{tabularx}{\linewidth}{>{\raggedright\arraybackslash}p{60mm}X}
\toprule
시험 (앞의 \code{test\_}는 생략) & 확인하는 것 \\
\midrule
\code{texts\_are\_enigma\_style} & 전문이 대문자뿐이고 \texttt{X}를 쓰고 움라우트가 없으며, 크립 후보가 실제로 전문 안에 있다 \\
\code{spec\_loops}, \code{default\_set\_...} & 루프 수의 계산, 기본 카드가 두 모양을 모두 덮는다 \\
\code{record\_structure\_\allowbreak and\_verification} & 카드마다: 모양, 단순 그래프, 넘어감 없음, 정답 정지 하나, 거짓 정지 2개 이하, \code{verify\_record} 통과 \\
\code{search\_is\_deterministic} & 카드마다: 씨앗과 시도 번호로 같은 카드가 다시 나온다 \\
\code{tampered\_record\_is\_rejected} & 고친 기록을 \code{verify\_record}가 거부한다 \\
\code{window\_letters\_\allowbreak are\_core\_plus\_ring} & 창 글자 = 코어 + 링, 유닛의 오른쪽 드럼이 크립 위치만큼 앞선다 \\
\code{solver\_makes\_the\_requested\_graph} & 플러그보드 탐색만으로 (선 6, 글자 4, 루프 3)이 나오고 실제 에니그마와 맞는다 \\
\code{card\_source\_\allowbreak has\_front\_back\_and\_answer} & 원고에 앞뒤 두 쪽, 그림, 유닛의 드럼 오프셋, 해독문이 있다 \\
\code{no\_personal\_paths\_in\_sources} & 소스와 기록에 개인 컴퓨터의 경로가 없다 \\
\bottomrule
\end{tabularx}
\end{center}

마지막 시험은 코드의 동작이 아니라 저장소의 위생을 본다. 생성기가 PDF를 만든 컴퓨터의 폴더 이름이 기록에 섞여
들어가면, 그 기록을 받은 다른 사람에게는 쓸모없는 정보가 개인 정보와 함께 퍼진다. 그래서 출력 폴더는 저장소 밖에 두고,
들어가서는 안 되는 문자열을 시험으로 막는다.

\begin{keybox}[title={이 장의 정리}]
\begin{itemize}
\item 카드는 메뉴의 모양(선, 글자, 루프)과 거짓 정지 수를 먼저 정하고, 그것을 주는 키를 찾는다.
\item 플러그보드는 $c_k = S(\sigma_k(S(p_k)))$를 거꾸로 풀어, 되추적으로 선 하나씩 정한다.
\item 찾은 키는 기준 코드(에니그마, 메뉴, 봄베, 체킹)로 처음부터 다시 계산해 확인한다.
\item 조각 카드는 느린 드럼을 정답 글자에 세워 676 위치만 돌고, 거짓 정지는 0--2개, 모두 체킹에서 걸러진다.
\item 씨앗과 시도 번호로 같은 카드를 다시 만들 수 있고, 기록은 생성기와 독립적으로 다시 확인된다.
\item 카드는 관람객의 교재이면서 세 기계의 시험지다.
\end{itemize}
\end{keybox}
\endgroup
